\documentclass[12pt]{article}
\usepackage[margin=3cm]{geometry}
\usepackage{amsmath, amssymb, amsthm}
\usepackage{tasks}
\usepackage[inline]{enumitem}
\usepackage{tcolorbox}
\tcbuselibrary{breakable}
\usepackage{fancyhdr}


% Initialize line spacing to 1.5
\linespread{1.5}

% Simple exercises environment using basic enumerate
\newcounter{exercisenum}
\newenvironment{exercises}{%
  \begin{tcolorbox}[colback=white,colframe=black,arc=0pt,breakable]
%   \noindent\textbf{Exercises}\par
  \begin{enumerate}[label=\arabic*),leftmargin=1cm,resume=exercises]
  \setcounter{enumi}{\value{exercisenum}}
}{%
  \setcounter{exercisenum}{\value{enumi}}
  \end{enumerate}
  \end{tcolorbox}
}

% Setup footer for first page only
\fancypagestyle{firstpage}{
  \fancyhf{}
  \fancyfoot[L]{\small MULTIPLES}
  \fancyfoot[R]{\small \quad Ewan Short \quad \today}
  \fancyfoot[C]{\small\thepage}
  \renewcommand{\headrulewidth}{0pt}
  \renewcommand{\footrulewidth}{0pt}
}

\begin{document}

\thispagestyle{firstpage}

\section{Whole Numbers and Multiplication}
\begin{itemize}
\item
A \underbar{whole number} is a counting number like $0,1,2,3,4$ etc. Numbers that can only be written as fractions or decimals are not whole numbers.
\item
\underbar{Multiplication} of whole numbers is repeated addition. For example,
\begin{equation*}
3\times 2 = 2 + 2 + 2 = 6, \quad\quad
4\times 1 = 1+1+1+1 = 4.
\end{equation*}
\end{itemize}

\begin{exercises}
\item
Circle the whole numbers.
$$4,\quad 739,\quad 2.123, \quad \frac{1}{3},\quad 8, \quad 911, \quad 0, \quad 0.01, \quad 1000000$$
\item
Write the following multiplications as repeated additions.
\begin{tasks}(2)
\task $2\times 4 = 4+4 = 8$
\task $2\times 2 = \phantom{2+2=4}$
\task $3\times 1 = \phantom{1+1+1=3}$
\task $5\times 2 = \phantom{2+2+2+2+2=10}$
\end{tasks}
\end{exercises}

\section{Multiples}
\begin{itemize}
\item
If we multiply a whole number by another whole number, we get a \underbar{multiple}. For example, the 5 times-table gives us some multiples of 5. 
\begin{align*}
5\times 1 &= 5, \quad \text{$5$ is a multiple of $5$ and of $1$,}\\
5\times 2 &= 10, \quad \text{$10$ is a multiple of $5$ and of $2$,}\\ 
\vdots\\
5 \times 12 &= 60, \quad \text{$60$ is a multiple of $5$ and of $12$.}
\end{align*}
\item
A \underbar{common multiple} of two numbers is a number that is a multiple of both. Below are some multiples of 5 and 10, with the common multiples in boxes.
\begin{align*}
\text{Multiples of 5:}\quad & 5,\quad\boxed{\boxed{10}},\quad 15,\quad\boxed{20},\quad 25,\quad\boxed{30},\quad 35,\quad\boxed{40},\quad 45,\quad\boxed{50},\ldots \\
\text{Multiples of 10:}\quad & \boxed{\boxed{10}},\quad \boxed{20},\quad \boxed{30},\quad \boxed{40},\quad \boxed{50},\quad \ldots
\end{align*}
\item
The \underbar{least common multiple} of two numbers is the smallest \underbar{positive} multiple of both. For example, the least common multiple of 5 and 10 is 10, and is double-boxed in the lists above.
\end{itemize}

\begin{exercises}
\item
\label{Ex:mult}
List the next 5 multiples of 4 and of 10. Your times-tables will help! 
\begin{align*}
\text{Multiples of 4:}\quad & 4, \quad 8,\quad 12,\quad 16,\quad 20,\quad \phantom{24, \quad 28, \quad 32, \quad 36, \quad 40}\\
\text{Multiples of 10:}\quad & 10,\quad 20,\quad 30,\quad 40,\quad 50, \phantom{60, \quad 70,\quad 80,\quad 90,\quad 100}
\end{align*}
\item
In your lists of multiples in \ref{Ex:mult} above, put a box around the common multiples, and a double box around the least common multiple.
\item
Find the least common multiples of the following numbers.
\begin{tasks}(4)
\task 2 and 3.
\task 4 and 6.
\task 4 and 8.
\task 3 and 7.
\end{tasks}
\end{exercises}

\end{document}